% Options for packages loaded elsewhere
% Options for packages loaded elsewhere
\PassOptionsToPackage{unicode}{hyperref}
\PassOptionsToPackage{hyphens}{url}
\PassOptionsToPackage{dvipsnames,svgnames,x11names}{xcolor}
%
\documentclass[
  spanish,
  10pt,
  letterpaper,
  twocolumn]{article}
\usepackage{xcolor}
\usepackage[top=1.0cm,bottom=1.5cm,left=1.5cm,right=1.5cm]{geometry}
\usepackage{amsmath,amssymb}
\setcounter{secnumdepth}{-\maxdimen} % remove section numbering
\usepackage{iftex}
\ifPDFTeX
  \usepackage[T1]{fontenc}
  \usepackage[utf8]{inputenc}
  \usepackage{textcomp} % provide euro and other symbols
\else % if luatex or xetex
  \usepackage{unicode-math} % this also loads fontspec
  \defaultfontfeatures{Scale=MatchLowercase}
  \defaultfontfeatures[\rmfamily]{Ligatures=TeX,Scale=1}
\fi
\usepackage{lmodern}
\ifPDFTeX\else
  % xetex/luatex font selection
  \setmainfont[]{Fira Sans}
\fi
% Use upquote if available, for straight quotes in verbatim environments
\IfFileExists{upquote.sty}{\usepackage{upquote}}{}
\IfFileExists{microtype.sty}{% use microtype if available
  \usepackage[]{microtype}
  \UseMicrotypeSet[protrusion]{basicmath} % disable protrusion for tt fonts
}{}
\usepackage{setspace}
\makeatletter
\@ifundefined{KOMAClassName}{% if non-KOMA class
  \IfFileExists{parskip.sty}{%
    \usepackage{parskip}
  }{% else
    \setlength{\parindent}{0pt}
    \setlength{\parskip}{6pt plus 2pt minus 1pt}}
}{% if KOMA class
  \KOMAoptions{parskip=half}}
\makeatother
% Make \paragraph and \subparagraph free-standing
\makeatletter
\ifx\paragraph\undefined\else
  \let\oldparagraph\paragraph
  \renewcommand{\paragraph}{
    \@ifstar
      \xxxParagraphStar
      \xxxParagraphNoStar
  }
  \newcommand{\xxxParagraphStar}[1]{\oldparagraph*{#1}\mbox{}}
  \newcommand{\xxxParagraphNoStar}[1]{\oldparagraph{#1}\mbox{}}
\fi
\ifx\subparagraph\undefined\else
  \let\oldsubparagraph\subparagraph
  \renewcommand{\subparagraph}{
    \@ifstar
      \xxxSubParagraphStar
      \xxxSubParagraphNoStar
  }
  \newcommand{\xxxSubParagraphStar}[1]{\oldsubparagraph*{#1}\mbox{}}
  \newcommand{\xxxSubParagraphNoStar}[1]{\oldsubparagraph{#1}\mbox{}}
\fi
\makeatother


\usepackage{longtable,booktabs,array}
\newcounter{none} % for unnumbered tables
\usepackage{calc} % for calculating minipage widths
% Correct order of tables after \paragraph or \subparagraph
\usepackage{etoolbox}
\makeatletter
\patchcmd\longtable{\par}{\if@noskipsec\mbox{}\fi\par}{}{}
\makeatother
% Allow footnotes in longtable head/foot
\IfFileExists{footnotehyper.sty}{\usepackage{footnotehyper}}{\usepackage{footnote}}
\makesavenoteenv{longtable}
\usepackage{graphicx}
\makeatletter
\newsavebox\pandoc@box
\newcommand*\pandocbounded[1]{% scales image to fit in text height/width
  \sbox\pandoc@box{#1}%
  \Gscale@div\@tempa{\textheight}{\dimexpr\ht\pandoc@box+\dp\pandoc@box\relax}%
  \Gscale@div\@tempb{\linewidth}{\wd\pandoc@box}%
  \ifdim\@tempb\p@<\@tempa\p@\let\@tempa\@tempb\fi% select the smaller of both
  \ifdim\@tempa\p@<\p@\scalebox{\@tempa}{\usebox\pandoc@box}%
  \else\usebox{\pandoc@box}%
  \fi%
}
% Set default figure placement to htbp
\def\fps@figure{htbp}
\makeatother



\ifLuaTeX
\usepackage[bidi=basic,shorthands=off]{babel}
\else
\usepackage[bidi=default,shorthands=off]{babel}
\fi
\ifPDFTeX
\else
\babelfont{rm}[]{Fira Sans}
\fi
\ifLuaTeX
  \usepackage{selnolig} % disable illegal ligatures
\fi


\setlength{\emergencystretch}{3em} % prevent overfull lines

\providecommand{\tightlist}{%
  \setlength{\itemsep}{0pt}\setlength{\parskip}{0pt}}



 


\usepackage{tikz}
\usepackage{graphicx}
\usepackage{amsmath}
\renewcommand{\arraystretch}{1.4}
\makeatletter
\@ifpackageloaded{caption}{}{\usepackage{caption}}
\AtBeginDocument{%
\ifdefined\contentsname
  \renewcommand*\contentsname{Tabla de contenidos}
\else
  \newcommand\contentsname{Tabla de contenidos}
\fi
\ifdefined\listfigurename
  \renewcommand*\listfigurename{Listado de Figuras}
\else
  \newcommand\listfigurename{Listado de Figuras}
\fi
\ifdefined\listtablename
  \renewcommand*\listtablename{Listado de Tablas}
\else
  \newcommand\listtablename{Listado de Tablas}
\fi
\ifdefined\figurename
  \renewcommand*\figurename{Figura}
\else
  \newcommand\figurename{Figura}
\fi
\ifdefined\tablename
  \renewcommand*\tablename{Tabla}
\else
  \newcommand\tablename{Tabla}
\fi
}
\@ifpackageloaded{float}{}{\usepackage{float}}
\floatstyle{ruled}
\@ifundefined{c@chapter}{\newfloat{codelisting}{h}{lop}}{\newfloat{codelisting}{h}{lop}[chapter]}
\floatname{codelisting}{Listado}
\newcommand*\listoflistings{\listof{codelisting}{Listado de Listados}}
\makeatother
\makeatletter
\makeatother
\makeatletter
\@ifpackageloaded{caption}{}{\usepackage{caption}}
\@ifpackageloaded{subcaption}{}{\usepackage{subcaption}}
\makeatother
\usepackage{bookmark}
\IfFileExists{xurl.sty}{\usepackage{xurl}}{} % add URL line breaks if available
\urlstyle{same}
% fallback for those not using the hyperref driver hyperxmp:
\makeatletter
\@ifundefined{xmpquote}{\newcommand{\xmpquote}[1]{#1}}{}
\makeatother
\hypersetup{
  pdftitle={Práctica -- Viabilidad de Semillas},
  pdflang={es},
  colorlinks=true,
  linkcolor={blue},
  filecolor={Maroon},
  citecolor={Blue},
  urlcolor={Blue},
  pdfcreator={LaTeX via pandoc}}


\title{\vspace{-1.5em}{\Huge \textbf{Práctica -- Viabilidad de
Semillas}}\vspace{-1.5em}}
\date{}
\begin{document}
\maketitle


\setstretch{1.2}
\textbf{Fecha}: \hrulefill \hfill \textbf{Grupo}: \hrulefill

\textbf{Nombre:} \hrulefill

\textbf{Objetivos:}

\begin{itemize}
\tightlist
\item
  Evaluar la viabilidad embrionaria en semillas de frijol
  (\emph{Phaseolus vulgaris}) o maíz (\emph{Zea mays}) mediante la
  prueba bioquímica del Cloruro de 2,3,5-trifeniltetrazolio (TZ).
\item
  Evaluar el efecto de la desnaturalización térmica enzimática sobre la
  capacidad respiratoria de las semillas.
\item
  Cuantificar el Porcentaje de Viabilidad (\(\%V\)) mediante estadística
  descriptiva (\(\bar{x} \pm s\)) y representaciones gráficas.
\end{itemize}

\subsection{1. Materiales}\label{materiales}

\begin{itemize}
\tightlist
\item
  Semillas de \emph{Phaseolus vulgaris} (frijol) o \emph{Zea mays}
  (maíz).
\item
  Kit de disección.
\item
  Guantes
\end{itemize}

\subsection{2. Aspectos Teóricos e
Introducción}\label{aspectos-teuxf3ricos-e-introducciuxf3n}

La prueba de \textbf{Tetrazolio} determina la viabilidad de las semillas
evaluando la actividad de las enzimas \textbf{deshidrogenasas} de la
cadena respiratoria mitocondrial. En células vivas e imbibidas, estas
enzimas transfieren protones (\(H^+\)) al Cloruro de
2,3,5-trifeniltetrazolio (incoloro y soluble), reduciéndolo a
\textbf{Formazán} (insoluble y de color rojo intenso). La exposición a
altas temperaturas causa la desnaturalización irreversible de la
estructura terciaria/cuaternaria de las proteínas y enzimas, anulando la
respiración celular y la reacción de tinción.

Defina brevemente:

• \textbf{Mecanismo Bioquímico de Reducción del Tetrazolio a Formazán:}

\hrulefill

\hrulefill

\hrulefill

\hrulefill

\hrulefill

\hrulefill

• \textbf{Efecto de la Temperatura sobre la Desnaturalización
Enzimática:}

\hrulefill

\hrulefill

\hrulefill

\hrulefill

\hrulefill

\hrulefill

• \textbf{Evaluación Topográfica (Diferencia entre daño en el eje
embrionario vs.~cotiledones):}

\hrulefill

\hrulefill

\hrulefill

\hrulefill

\hrulefill

\hrulefill

\subsection{3. Procedimiento
Experimental}\label{procedimiento-experimental}

\subsubsection{A. Imbibición y
Preacondicionamiento}\label{a.-imbibiciuxf3n-y-preacondicionamiento}

\begin{enumerate}
\def\labelenumi{\arabic{enumi}.}
\tightlist
\item
  Coloque 60 semillas en agua a temperatura ambiente durante 18--24
  horas antes de la práctica para activar el metabolismo celular y
  facilitar el corte.
\end{enumerate}

\subsubsection{B. Aplicación del Tratamiento
Térmico}\label{b.-aplicaciuxf3n-del-tratamiento-tuxe9rmico}

\begin{enumerate}
\def\labelenumi{\arabic{enumi}.}
\tightlist
\item
  \textbf{Grupo Tratamiento Térmico:} Sumerja 30 semillas imbibidas en
  un vaso de precipitado con agua caliente a \(80-90^\circ\text{C}\) en
  baño María durante \textbf{15 minutos} para inducir la inactivación
  térmica de las deshidrogenasas.
\item
  Retire las semillas del agua caliente y déjelas enfriar a temperatura
  ambiente en agua destilada.
\end{enumerate}

\subsubsection{C. Corte e Incubación en Solución de
TZ}\label{c.-corte-e-incubaciuxf3n-en-soluciuxf3n-de-tz}

\begin{enumerate}
\def\labelenumi{\arabic{enumi}.}
\tightlist
\item
  Con el bisturí, realice un corte longitudinal sagital exponiendo
  claramente el eje embrionario y los cotiledones.
\item
  Coloque la mitad de cada semilla que contenga el embrión en cajas de
  Petri cubiertas con 5 mL de solución de TZ al 1\% e incube a
  \(35^\circ\text{C}\) en oscuridad durante 60 minutos:

  \begin{itemize}
  \tightlist
  \item
    \textbf{T1 (Control - Sin Calor):} 3 cajas de Petri de 10 semillas
    cada una.
  \item
    \textbf{T2 (Tratamiento Térmico - Con Calor):} 3 cajas de Petri de
    10 semillas cada una.
  \end{itemize}
\end{enumerate}

\begin{center}
\begin{tikzpicture}[scale=0.78]
  % SEMILLA VIABLE (Control - Tinción Completa)
  \begin{scope}[shift={(0,0)}]
    \draw[thick, fill=red!70!black] (0,0) .. controls (-0.6,0.8) and (-0.6,1.8) .. (0,2.2)
                                     .. controls (0.6,1.8) and (0.6,0.8) .. (0,0) -- cycle;
    \draw[thick, fill=red!90!black] (-0.05,0.3) .. controls (-0.2,0.8) .. (-0.05,1.2) -- (0.05,1.2) .. controls (0.2,0.8) .. (0.05,0.3) -- cycle;
    
    \node[font=\fontsize{6}{7}\bfseries\selectfont, align=center] at (0, -0.4) {Control (Sin Calor)};
    \node[font=\fontsize{5}{6}\selectfont, gray!50!black, align=center] at (0, -0.7) {Tejido Vivo / Respirando};
  \end{scope}

  % SEMILLA PARCIAL / DAÑO MODERADO
  \begin{scope}[shift={(3,0)}]
    \draw[thick, fill=gray!20] (0,0) .. controls (-0.6,0.8) and (-0.6,1.8) .. (0,2.2)
                                   .. controls (0.6,1.8) and (0.6,0.8) .. (0,0) -- cycle;
    \fill[red!60!black] (-0.02,0.02) .. controls (-0.55,0.8) and (-0.55,1.4) .. (-0.02,1.8) -- cycle;
    \draw[thick, fill=red!80!black] (-0.05,0.3) .. controls (-0.2,0.8) .. (-0.05,1.2) -- (0.05,1.2) .. controls (0.2,0.8) .. (0.05,0.3) -- cycle;

    \node[font=\fontsize{6}{7}\bfseries\selectfont, align=center] at (0, -0.4) {Calor Parcial};
    \node[font=\fontsize{5}{6}\selectfont, gray!50!black, align=center] at (0, -0.7) {Desnaturalización Incompleta};
  \end{scope}

  % SEMILLA NO VIABLE (Tratamiento Térmico - Sin tinción)
  \begin{scope}[shift={(6,0)}]
    \draw[thick, fill=gray!15] (0,0) .. controls (-0.6,0.8) and (-0.6,1.8) .. (0,2.2)
                                     .. controls (0.6,1.8) and (0.6,0.8) .. (0,0) -- cycle;
    \draw[thick, fill=gray!30] (-0.05,0.3) .. controls (-0.2,0.8) .. (-0.05,1.2) -- (0.05,1.2) .. controls (0.2,0.8) .. (0.05,0.3) -- cycle;

    \node[font=\fontsize{6}{7}\bfseries\selectfont, align=center] at (0, -0.4) {Tratado con Calor};
    \node[font=\fontsize{5}{6}\selectfont, gray!50!black, align=center] at (0, -0.7) {Enzimas Inactivadas};
  \end{scope}
\end{tikzpicture}
\end{center}

\subsubsection{D. Evaluación
Topográfica}\label{d.-evaluaciuxf3n-topogruxe1fica}

\begin{enumerate}
\def\labelenumi{\arabic{enumi}.}
\tightlist
\item
  Enjuague las semillas con agua destilada y examine bajo el
  estereoscopio.
\item
  Clasifique cada semilla como \textbf{Viable} (tinción roja intensa en
  el eje embrionario) o \textbf{No Viable} (sin tinción o tejido
  blanquecino/necrosado por efecto térmico).
\item
  Calcule el porcentaje de viabilidad (\(\%V\)) por réplica (caja de
  Petri) y determine el promedio (\(\bar{x}\)) y la desviación estándar
  (\(s\)) para cada tratamiento:
\end{enumerate}

\[\%V = \frac{\text{Semillas Viables}}{\text{Total Semillas}} \times 100\]

\[\bar{x} = \frac{\sum_{i=1}^{n} x_i}{n} \qquad s = \sqrt{\frac{\sum_{i=1}^{n} (x_i - \bar{x})^2}{n - 1}}\]

\subsection{4. Registro de Resultados}\label{registro-de-resultados}

\subsubsection{Tabla 1. Porcentaje de Viabilidad (\%V) por
Tratamiento}\label{tabla-1.-porcentaje-de-viabilidad-v-por-tratamiento}

\begin{center}
\resizebox{\linewidth}{!}{%
\begin{tabular}{l|l|p{2em}|p{2em}|p{2em}|p{2em}|p{2em}}
\hline
 &  & \multicolumn{3}{c|}{\textbf{Réplicas}} & \textbf{Prom} & \textbf{DE} \\
\textbf{Tratamiento} & \textbf{Categoría} & \textbf{A} & \textbf{B} & \textbf{C} & \textbf{$\bar{x}$} & \textbf{$s$} \\
\hline
1. Control & Viables (\%) & & & & & \\[1ex]
\cline{2-7}
 & No Viables (\%) & & & & & \\[1ex]
\hline
2. Calor & Viables (\%) & & & & & \\[1ex]
\cline{2-7}
 & No Viables (\%) & & & & & \\[1ex]
\hline
\end{tabular}%
}
\end{center}

Interprete los resultados del porcentaje de semillas viables según
tratamiento:

\hrulefill

\hrulefill

\hrulefill

\hrulefill

\hrulefill

\subsubsection{Tabla 2. Descripción Topográfica de Patrones de Tinción y
Coloración}\label{tabla-2.-descripciuxf3n-topogruxe1fica-de-patrones-de-tinciuxf3n-y-coloraciuxf3n}

\begin{center}
\resizebox{\linewidth}{!}{%
\begin{tabular}{l|p{14em}}
\hline
\textbf{Tratamiento} & \textbf{Observación Microscópica y Localización del Formazán} \\
\hline
1. Control (Sin Calor) & \\[40pt]
\hline
2. Tratamiento Térmico & \\[40pt]
\hline
\end{tabular}%
}
\end{center}

Interprete los resultados de la observación microscópica y localización
del Formazán:

\hrulefill

\hrulefill

\hrulefill

\hrulefill

\hrulefill

\subsection{5. Análisis Gráfico}\label{anuxe1lisis-gruxe1fico}

Represente la viabilidad de las semillas utilizando puntos (promedio) y
barras de error (desviación estándar) por tratamiento.

\subsubsection{Gráfica 1: Porcentaje de Viabilidad (\%V) por
Tratamiento}\label{gruxe1fica-1-porcentaje-de-viabilidad-v-por-tratamiento}

\begin{center}
\begin{tikzpicture}[scale=0.85]
  % Grid ggplot2
  \draw[gray!20, step=0.7cm] (0,0) grid (7,5);
  \draw[gray!50, step=1.4cm] (0,0) grid (7,5);
  \draw[thick] (0,0) rectangle (7,5);
  
  % Eje X
  \node[below, font=\small] at (2.1,0) {\textbf{Control}};
  \node[below, font=\small] at (4.9,0) {\textbf{Calor}};
  \node[below=12pt, font=\small] at (3.25,0) {\textbf{Tratamiento Experimental}};
  
  % Eje Y
  \node[left, font=\tiny] at (0,0) {0};
  \node[left, font=\tiny] at (0,1) {25};
  \node[left, font=\tiny] at (0,2) {50};
  \node[left, font=\tiny] at (0,3) {75};
  \node[left, font=\tiny] at (0,4) {100};
  \node[rotate=90, above=10pt, font=\small] at (0,2.0) {\textbf{Viabilidad Seminal (\%)}};
\end{tikzpicture}
\end{center}

Interprete los resultados gráficos:

\hrulefill

\hrulefill

\hrulefill

\hrulefill

\subsection{6. Actividades
Complementarias}\label{actividades-complementarias}

\begin{enumerate}
\def\labelenumi{\arabic{enumi}.}
\tightlist
\item
  Explique a nivel molecular por qué el choque térmico
  (\(80-90^\circ\text{C}\)) impide la reducción del Cloruro de
  Tetrazolio a Formazán.
\end{enumerate}

\hrulefill

\hrulefill

\hrulefill

\hrulefill

\hrulefill

\hrulefill

\begin{enumerate}
\def\labelenumi{\arabic{enumi}.}
\setcounter{enumi}{1}
\tightlist
\item
  ¿Qué diferencia existiría en el patrón de tinción si las semillas se
  someten a un calor moderado (\(50^\circ\text{C}\)) en lugar de un
  calor extremo (\(90^\circ\text{C}\))?
\end{enumerate}

\hrulefill

\hrulefill

\hrulefill

\hrulefill

\hrulefill

\hrulefill

\begin{enumerate}
\def\labelenumi{\arabic{enumi}.}
\setcounter{enumi}{2}
\tightlist
\item
  ¿Por qué la prueba de Tetrazolio permite diagnosticar la muerte
  celular antes de que se observen cambios morfológicos externos en la
  semilla?
\end{enumerate}

\hrulefill

\hrulefill

\hrulefill

\hrulefill

\hrulefill

\hrulefill

\subsection{7. Conclusión}\label{conclusiuxf3n}

Interprete el impacto de la desnaturalización térmica sobre la actividad
enzimática respiratoria y analice la confiabilidad de la prueba de TZ
para diferenciar semillas viables de no viables:

\hrulefill

\hrulefill

\hrulefill

\hrulefill

\hrulefill

\hrulefill

\subsection{8. Bibliografía}\label{bibliografuxeda}

\begin{itemize}
\tightlist
\item
  Azcón-Bieto, J., \& Talón, M. (2008). \emph{Fundamentos de Fisiología
  Vegetal} (2.ª ed.). McGraw-Hill.
\item
  Franca-Neto, J. B., \& Krzyzanowski, F. C. (2019). \emph{Tetrazolium
  test for seed viability and vigor}. Embrapa Soja.
\item
  ISTA. (2020). \emph{International Rules for Seed Testing}.
  International Seed Testing Association.
\item
  Taiz, L., \& Zeiger, E. (2002). \emph{Plant Physiology} (3rd ed.).
  Sinauer Associates.
\end{itemize}




\end{document}
